\documentclass[11pt]{article}
\usepackage[utf8]{inputenc}
\usepackage{amsmath, amssymb, geometry, graphicx, enumitem}
\geometry{a4paper, margin=1in}

\title{Math Problem Set: Triple Integrals in Spherical Coordinates}
\author{}
\date{}

\begin{document}

\maketitle

\section*{Concept Checklist (Start)}
The following concepts and problem types extracted from the source material are tested in this problem set:
\begin{itemize}
    \item \textbf{C1:} Point Conversion (Spherical to Cartesian)
    \item \textbf{C2:} Point Conversion (Cartesian to Spherical)
    \item \textbf{C3:} Identifying Surfaces from Spherical Equations
    \item \textbf{C4:} Converting Cartesian Equations to Spherical Equations
    \item \textbf{C5:} Geometric Visualization \& Setting Up Inequalities for 3D Regions
    \item \textbf{C6:} Evaluating Separable Triple Integrals (Constant Bounds / Spherical Rectangles)
    \item \textbf{C7:} Evaluating Triple Integrals over Spheres and Balls
    \item \textbf{C8:} Evaluating Triple Integrals bounded by Cones and Spheres
    \item \textbf{C9:} Converting Cartesian Triple Integrals to Spherical Integrals
    \item \textbf{C10:} Advanced Integration Techniques (Integration by Parts, Odd Powers of Sine/Cosine)
    \item \textbf{C11:} Real-World Applications \& Modeling (Probability \& Expected Distance)
    \item \textbf{C12:} Center of Mass / Centroids
    \item \textbf{C13:} "Find the Flaw" / Diagnostic Testing
\end{itemize}

\hrule \vspace{0.5cm}

\section*{The 30 Problem Sets}

\begin{enumerate}[label=\textbf{\arabic*.}]

% ----- CONCEPT 1 -----
\item Plot the point whose spherical coordinates are given and find its rectangular coordinates: $(4, \frac{\pi}{3}, \frac{\pi}{6})$.\\
\textbf{Solution:} We are given $(\rho, \theta, \phi) = (4, \frac{\pi}{3}, \frac{\pi}{6})$. Apply the conversion formulas:
\begin{align*}
x &= \rho \sin\phi \cos\theta = 4 \sin\left(\frac{\pi}{6}\right) \cos\left(\frac{\pi}{3}\right) = 4 \left(\frac{1}{2}\right) \left(\frac{1}{2}\right) = 1 \\
y &= \rho \sin\phi \sin\theta = 4 \sin\left(\frac{\pi}{6}\right) \sin\left(\frac{\pi}{3}\right) = 4 \left(\frac{1}{2}\right) \left(\frac{\sqrt{3}}{2}\right) = \sqrt{3} \\
z &= \rho \cos\phi = 4 \cos\left(\frac{\pi}{6}\right) = 4 \left(\frac{\sqrt{3}}{2}\right) = 2\sqrt{3}
\end{align*}
\textbf{Final Answer:} $(1, \sqrt{3}, 2\sqrt{3})$

\item Plot the point whose spherical coordinates are given and find its rectangular coordinates: $(2, \frac{3\pi}{4}, \frac{3\pi}{4})$.\\
\textbf{Solution:} We are given $(\rho, \theta, \phi) = (2, \frac{3\pi}{4}, \frac{3\pi}{4})$.
\begin{align*}
x &= \rho \sin\phi \cos\theta = 2 \sin\left(\frac{3\pi}{4}\right) \cos\left(\frac{3\pi}{4}\right) = 2 \left(\frac{\sqrt{2}}{2}\right) \left(-\frac{\sqrt{2}}{2}\right) = -1 \\
y &= \rho \sin\phi \sin\theta = 2 \sin\left(\frac{3\pi}{4}\right) \sin\left(\frac{3\pi}{4}\right) = 2 \left(\frac{\sqrt{2}}{2}\right) \left(\frac{\sqrt{2}}{2}\right) = 1 \\
z &= \rho \cos\phi = 2 \cos\left(\frac{3\pi}{4}\right) = 2 \left(-\frac{\sqrt{2}}{2}\right) = -\sqrt{2}
\end{align*}
\textbf{Final Answer:} $(-1, 1, -\sqrt{2})$

% ----- CONCEPT 2 -----
\item Change from rectangular to spherical coordinates: $(0, 3, 3)$. Let $\rho \geq 0, 0 \le \theta \le 2\pi, 0 \le \phi \le \pi$.\\
\textbf{Solution:} Given $(x, y, z) = (0, 3, 3)$.
\begin{align*}
\rho &= \sqrt{x^2+y^2+z^2} = \sqrt{0 + 9 + 9} = \sqrt{18} = 3\sqrt{2} \\
\cos\phi &= \frac{z}{\rho} = \frac{3}{3\sqrt{2}} = \frac{1}{\sqrt{2}} \implies \phi = \frac{\pi}{4} \\
\tan\theta &= \frac{y}{x} = \frac{3}{0} \text{ (undefined)}
\end{align*}
Since $x=0$ and $y>0$, the point lies on the positive y-axis in the $xy$-plane, meaning $\theta = \frac{\pi}{2}$.\\
\textbf{Final Answer:} $(3\sqrt{2}, \frac{\pi}{2}, \frac{\pi}{4})$

\item Change from rectangular to spherical coordinates: $(-2, 2, 2\sqrt{6})$.\\
\textbf{Solution:} Given $(x, y, z) = (-2, 2, 2\sqrt{6})$.
\begin{align*}
\rho &= \sqrt{(-2)^2 + (2)^2 + (2\sqrt{6})^2} = \sqrt{4 + 4 + 24} = \sqrt{32} = 4\sqrt{2} \\
\cos\phi &= \frac{z}{\rho} = \frac{2\sqrt{6}}{4\sqrt{2}} = \frac{\sqrt{3}}{2} \implies \phi = \frac{\pi}{6} \\
\tan\theta &= \frac{y}{x} = \frac{2}{-2} = -1
\end{align*}
Since $x<0$ and $y>0$, $\theta$ is in Quadrant II, so $\theta = \frac{3\pi}{4}$.\\
\textbf{Final Answer:} $(4\sqrt{2}, \frac{3\pi}{4}, \frac{\pi}{6})$

% ----- CONCEPT 3 -----
\item Identify the surface whose equation is given: $\rho \sin\phi = 4$.\\
\textbf{Solution:} In cylindrical/spherical relations, the radial distance from the z-axis is $r = \rho \sin\phi$. Substitute this in:
$r = 4 \implies \sqrt{x^2+y^2} = 4 \implies x^2 + y^2 = 16$.
This represents a cylinder of radius 4 centered along the z-axis.\\
\textbf{Final Answer:} The cylinder $x^2 + y^2 = 16$.

\item Identify the surface whose equation is given: $\rho = 2 \cos\phi$.\\
\textbf{Solution:} Multiply both sides by $\rho$:
$\rho^2 = 2\rho\cos\phi$.
Substitute spherical conversions $\rho^2 = x^2+y^2+z^2$ and $z = \rho\cos\phi$:
$x^2+y^2+z^2 = 2z \implies x^2+y^2+z^2-2z = 0$.
Complete the square for $z$: $x^2+y^2+(z-1)^2 = 1$.
This is a sphere of radius 1 centered at $(0,0,1)$.\\
\textbf{Final Answer:} The sphere $x^2+y^2+(z-1)^2=1$.

% ----- CONCEPT 4 -----
\item Write the equation in spherical coordinates: $x^2+y^2+z^2 = 9z$.\\
\textbf{Solution:} Use the identities $\rho^2 = x^2+y^2+z^2$ and $z = \rho\cos\phi$:
$\rho^2 = 9\rho\cos\phi$.
Assuming $\rho \neq 0$, divide by $\rho$:
$\rho = 9\cos\phi$.\\
\textbf{Final Answer:} $\rho = 9\cos\phi$

\item Write the equation in spherical coordinates: $z = 3x^2 + 3y^2$.\\
\textbf{Solution:} Factor the right side: $z = 3(x^2+y^2)$.
Substitute $z = \rho\cos\phi$ and $x^2+y^2 = (\rho\sin\phi)^2$:
$\rho\cos\phi = 3\rho^2\sin^2\phi$.
Divide by $\rho$:
$\cos\phi = 3\rho\sin^2\phi \implies \rho = \frac{\cos\phi}{3\sin^2\phi} = \frac{1}{3}\csc\phi\cot\phi$.\\
\textbf{Final Answer:} $\rho = \frac{1}{3}\csc\phi\cot\phi$

% ----- CONCEPT 5 -----
\item A solid lies between the spheres $\rho=2$ and $\rho=5$, and above the cone $z = \sqrt{x^2+y^2}$. Write a description of the solid using inequalities in spherical coordinates.\\
\textbf{Solution:} The spheres limit $\rho$ directly: $2 \le \rho \le 5$.
Convert the cone $z = \sqrt{x^2+y^2}$:
$\rho\cos\phi = \sqrt{\rho^2\sin^2\phi} = \rho\sin\phi \implies 1 = \tan\phi \implies \phi = \frac{\pi}{4}$.
"Above" the cone means it is closer to the positive z-axis, meaning angles smaller than $\pi/4$. Thus, $0 \le \phi \le \frac{\pi}{4}$. No $\theta$ restrictions are given, so $0 \le \theta \le 2\pi$.\\
\textbf{Final Answer:} $2 \le \rho \le 5$, $0 \le \phi \le \frac{\pi}{4}$, $0 \le \theta \le 2\pi$.

\item Sketch or describe the solid described by the inequalities: $1 \le \rho \le 3$, $0 \le \phi \le \frac{\pi}{3}$, $0 \le \theta \le \pi$.\\
\textbf{Solution:}
- $0 \le \theta \le \pi$: The solid is restricted to the half-space where $y \ge 0$ (Quadrants I and II of the xy-plane).
- $0 \le \phi \le \pi/3$: The solid lies within a cone opening upwards from the origin with an angle of $\pi/3$ from the z-axis.
- $1 \le \rho \le 3$: The solid is "hollowed out", lying between a smaller sphere of radius 1 and a larger sphere of radius 3. \\
\textbf{Final Answer:} A wedge of a spherical cone, located in the $y \ge 0$ half-space, bounded between spheres of radius 1 and 3.

% ----- CONCEPT 6 -----
\item Evaluate the integral: $\int_0^{\pi} \int_0^{\pi/2} \int_0^2 \rho^3 \sin\phi \cos\phi \, d\rho \, d\phi \, d\theta$.\\
\textbf{Solution:} Because all limits are constants and the integrand factors, we separate the integrals:
$= \left( \int_0^\pi d\theta \right) \left( \int_0^{\pi/2} \sin\phi \cos\phi \, d\phi \right) \left( \int_0^2 \rho^3 \, d\rho \right)$
1. $\int_0^\pi d\theta = \pi$
2. Let $u = \sin\phi, du = \cos\phi \, d\phi$. Bounds $0 \to 1$. $\int_0^1 u \, du = [\frac{u^2}{2}]_0^1 = \frac{1}{2}$.
3. $\int_0^2 \rho^3 \, d\rho = [\frac{\rho^4}{4}]_0^2 = \frac{16}{4} = 4$.
Multiply: $(\pi)(\frac{1}{2})(4) = 2\pi$.\\
\textbf{Final Answer:} $2\pi$

\item Evaluate the integral: $\int_0^{2\pi} \int_0^{\pi/4} \int_0^3 \rho^2 \sin\phi \, d\rho \, d\phi \, d\theta$.\\
\textbf{Solution:} Split the separable integral:
$= \left( \int_0^{2\pi} d\theta \right) \left( \int_0^{\pi/4} \sin\phi \, d\phi \right) \left( \int_0^3 \rho^2 \, d\rho \right)$
1. $\int_0^{2\pi} d\theta = 2\pi$.
2. $\int_0^{\pi/4} \sin\phi \, d\phi = [-\cos\phi]_0^{\pi/4} = -\frac{\sqrt{2}}{2} - (-1) = 1 - \frac{\sqrt{2}}{2}$.
3. $\int_0^3 \rho^2 \, d\rho = [\frac{\rho^3}{3}]_0^3 = \frac{27}{3} = 9$.
Multiply: $2\pi \left(1 - \frac{\sqrt{2}}{2}\right) (9) = 18\pi(1 - \frac{\sqrt{2}}{2}) = 9\pi(2-\sqrt{2})$.\\
\textbf{Final Answer:} $9\pi(2-\sqrt{2})$

\item Evaluate the integral: $\int_0^{\pi/2} \int_{\pi/4}^{\pi/2} \int_1^2 \rho^4 \sin^3\phi \, d\rho \, d\phi \, d\theta$.\\
\textbf{Solution:} Split into three integrals:
1. $\int_0^{\pi/2} d\theta = \frac{\pi}{2}$.
2. $\int_1^2 \rho^4 d\rho = [\frac{\rho^5}{5}]_1^2 = \frac{32-1}{5} = \frac{31}{5}$.
3. $\int_{\pi/4}^{\pi/2} \sin^3\phi d\phi$. Use the odd power trick: $\sin^3\phi = (1-\cos^2\phi)\sin\phi$. Let $u=\cos\phi$, $du=-\sin\phi d\phi$. Limits: $\cos(\pi/4) = \frac{\sqrt{2}}{2}$ to $\cos(\pi/2) = 0$.
$\int_{\sqrt{2}/2}^0 (1-u^2)(-du) = \int_0^{\sqrt{2}/2} (1-u^2)du = [u - \frac{u^3}{3}]_0^{\sqrt{2}/2} = \frac{\sqrt{2}}{2} - \frac{2\sqrt{2}/8}{3} = \frac{\sqrt{2}}{2} - \frac{\sqrt{2}}{12} = \frac{5\sqrt{2}}{12}$.
Multiply: $(\frac{\pi}{2}) (\frac{31}{5}) (\frac{5\sqrt{2}}{12}) = \frac{31\pi\sqrt{2}}{24}$.\\
\textbf{Final Answer:} $\frac{31\pi\sqrt{2}}{24}$

% ----- CONCEPT 7 & 8 -----
\item Find the volume of the solid that lies within the sphere $x^2+y^2+z^2=25$ and outside the cone $z=\sqrt{x^2+y^2}$. (Assume $z \ge 0$).\\
\textbf{Solution:} 
Bounds: Sphere $\rho = 5$. Cone $\phi = \pi/4$. "Outside" the cone but in the upper half-space means dropping further down toward the xy-plane: $\pi/4 \le \phi \le \pi/2$. Full rotation: $0 \le \theta \le 2\pi$. Volume element $dV = \rho^2\sin\phi d\rho d\phi d\theta$.
$V = \int_0^{2\pi} d\theta \int_{\pi/4}^{\pi/2} \sin\phi d\phi \int_0^5 \rho^2 d\rho$
$V = (2\pi) [-\cos\phi]_{\pi/4}^{\pi/2} [\frac{125}{3}] = 2\pi (0 - (-\frac{\sqrt{2}}{2})) (\frac{125}{3}) = 2\pi (\frac{\sqrt{2}}{2}) (\frac{125}{3}) = \frac{125\pi\sqrt{2}}{3}$.\\
\textbf{Final Answer:} $\frac{125\pi\sqrt{2}}{3}$

\item Evaluate $\iiint_B (x^2+y^2+z^2)^{3/2} dV$, where $B$ is the unit ball.\\
\textbf{Solution:} 
Bounds for unit ball: $0 \le \rho \le 1, 0 \le \phi \le \pi, 0 \le \theta \le 2\pi$.
Integrand: $(x^2+y^2+z^2)^{3/2} = (\rho^2)^{3/2} = \rho^3$.
Include Jacobian $dV = \rho^2\sin\phi \,d\rho \,d\phi \,d\theta$.
$I = \int_0^{2\pi} \int_0^\pi \int_0^1 (\rho^3)(\rho^2\sin\phi) d\rho \,d\phi \,d\theta = \int_0^{2\pi} d\theta \int_0^\pi \sin\phi d\phi \int_0^1 \rho^5 d\rho$
$I = (2\pi) [-\cos\phi]_0^\pi [\frac{\rho^6}{6}]_0^1 = (2\pi) (2) (\frac{1}{6}) = \frac{4\pi}{6} = \frac{2\pi}{3}$.\\
\textbf{Final Answer:} $\frac{2\pi}{3}$

\item Evaluate $\iiint_E z \, dV$, where $E$ is the upper hemisphere of radius 4.\\
\textbf{Solution:} 
Bounds for upper hemisphere: $0 \le \rho \le 4, 0 \le \phi \le \pi/2, 0 \le \theta \le 2\pi$.
Integrand: $z = \rho\cos\phi$. Multiply by $dV = \rho^2\sin\phi \,d\rho \,d\phi \,d\theta$.
$I = \int_0^{2\pi} \int_0^{\pi/2} \int_0^4 (\rho\cos\phi)(\rho^2\sin\phi) d\rho \,d\phi \,d\theta = \int_0^{2\pi} d\theta \int_0^{\pi/2} \sin\phi\cos\phi d\phi \int_0^4 \rho^3 d\rho$
$I = (2\pi) [\frac{\sin^2\phi}{2}]_0^{\pi/2} [\frac{\rho^4}{4}]_0^4 = (2\pi) (\frac{1}{2}) (\frac{256}{4}) = \pi(64) = 64\pi$.\\
\textbf{Final Answer:} $64\pi$

\item Evaluate $\iiint_E (x^2+y^2) dV$, where $E$ lies between the spheres $\rho=1$ and $\rho=2$.\\
\textbf{Solution:}
Bounds: $1 \le \rho \le 2, 0 \le \phi \le \pi, 0 \le \theta \le 2\pi$.
Integrand: $x^2+y^2 = (\rho\sin\phi\cos\theta)^2 + (\rho\sin\phi\sin\theta)^2 = \rho^2\sin^2\phi$.
Multiply by $dV = \rho^2\sin\phi \,d\rho \,d\phi \,d\theta$:
$I = \int_0^{2\pi} d\theta \int_0^\pi \sin^3\phi d\phi \int_1^2 \rho^4 d\rho$
1. $\int_0^{2\pi} d\theta = 2\pi$.
2. $\int_0^\pi \sin^3\phi d\phi = \int_0^\pi (1-\cos^2\phi)\sin\phi d\phi = [-\cos\phi + \frac{\cos^3\phi}{3}]_0^\pi = (1 - \frac{1}{3}) - (-1 + \frac{1}{3}) = \frac{4}{3}$.
3. $\int_1^2 \rho^4 d\rho = [\frac{\rho^5}{5}]_1^2 = \frac{32-1}{5} = \frac{31}{5}$.
Multiply: $(2\pi)(\frac{4}{3})(\frac{31}{5}) = \frac{248\pi}{15}$.\\
\textbf{Final Answer:} $\frac{248\pi}{15}$

% ----- CONCEPT 9 -----
\item Evaluate the integral by changing to spherical coordinates: $\int_{-3}^3 \int_{-\sqrt{9-x^2}}^{\sqrt{9-x^2}} \int_0^{\sqrt{9-x^2-y^2}} z \sqrt{x^2+y^2+z^2} \, dz \, dy \, dx$.\\
\textbf{Solution:}
Step 1: Region. The x and y bounds describe a full circle of radius 3 in the xy-plane ($0 \le \theta \le 2\pi$). The z bounds go from $0$ to the upper hemisphere of radius 3 ($0 \le \phi \le \pi/2$). Spherical radius $0 \le \rho \le 3$.
Step 2: Integrand. $z\sqrt{x^2+y^2+z^2} = (\rho\cos\phi)(\rho) = \rho^2\cos\phi$. Multiply by Jacobian $\rho^2\sin\phi \,d\rho \,d\phi \,d\theta$.
$I = \int_0^{2\pi} d\theta \int_0^{\pi/2} \sin\phi\cos\phi d\phi \int_0^3 \rho^4 d\rho$
$I = (2\pi) [\frac{\sin^2\phi}{2}]_0^{\pi/2} [\frac{\rho^5}{5}]_0^3 = (2\pi) (\frac{1}{2}) (\frac{243}{5}) = \frac{243\pi}{5}$.\\
\textbf{Final Answer:} $\frac{243\pi}{5}$

\item Convert the Cartesian integral to spherical coordinates and evaluate: 
$\int_0^2 \int_0^{\sqrt{4-x^2}} \int_{\sqrt{x^2+y^2}}^{\sqrt{8-x^2-y^2}} (x^2+y^2+z^2) \, dz \, dy \, dx$.\\
\textbf{Solution:}
Step 1: Region. $0 \le x \le 2, 0 \le y \le \sqrt{4-x^2}$ defines the first quadrant of a circle of radius 2 ($0 \le \theta \le \pi/2$). 
$z$ is bounded below by cone $z=\sqrt{x^2+y^2} \implies \phi=\pi/4$. 
$z$ bounded above by sphere $z=\sqrt{8-x^2-y^2} \implies \rho=\sqrt{8}=2\sqrt{2}$. 
Notice the intersection of cone and sphere is $2z^2 = 8 \implies z=2$, mapping exactly to radius $r=2$ in the xy-plane, justifying our xy bounds. Thus $0 \le \phi \le \pi/4$ and $0 \le \rho \le 2\sqrt{2}$.
Step 2: Integrand $(x^2+y^2+z^2) = \rho^2$. Multiply by $dV$.
$I = \int_0^{\pi/2} d\theta \int_0^{\pi/4} \sin\phi d\phi \int_0^{2\sqrt{2}} \rho^4 d\rho$
$I = (\frac{\pi}{2}) [-\cos\phi]_0^{\pi/4} [\frac{\rho^5}{5}]_0^{2\sqrt{2}} = (\frac{\pi}{2}) (1 - \frac{\sqrt{2}}{2}) (\frac{128\sqrt{2}}{5}) = \frac{\pi(2-\sqrt{2})}{4} \cdot \frac{128\sqrt{2}}{5} = \frac{32\pi(2\sqrt{2}-2)}{5}$.\\
\textbf{Final Answer:} $\frac{64\pi(\sqrt{2}-1)}{5}$

\item Convert to spherical coordinates and evaluate: $\int_{-1}^1 \int_{-\sqrt{1-y^2}}^{\sqrt{1-y^2}} \int_{-\sqrt{1-x^2-y^2}}^0 e^{(x^2+y^2+z^2)^{3/2}} \, dz \, dx \, dy$.\\
\textbf{Solution:}
Note the $dx dy$ order. $y \in [-1, 1], x \in [-\sqrt{1-y^2}, \sqrt{1-y^2}]$ forms a full circle of radius 1 ($0 \le \theta \le 2\pi$). 
$z \in [-\sqrt{1-x^2-y^2}, 0]$ defines the lower unit hemisphere ($\pi/2 \le \phi \le \pi$, and $0 \le \rho \le 1$).
Integrand: $e^{(\rho^2)^{3/2}} = e^{\rho^3}$. Multiply by $dV$.
$I = \int_0^{2\pi} d\theta \int_{\pi/2}^\pi \sin\phi d\phi \int_0^1 e^{\rho^3} \rho^2 d\rho$
1. $\int_0^{2\pi} d\theta = 2\pi$.
2. $\int_{\pi/2}^\pi \sin\phi d\phi = [-\cos\phi]_{\pi/2}^\pi = -(-1) - 0 = 1$.
3. $\int_0^1 e^{\rho^3} \rho^2 d\rho$. Let $u=\rho^3, du=3\rho^2 d\rho$. $\int_0^1 \frac{1}{3}e^u du = \frac{e-1}{3}$.
Multiply: $2\pi(1)(\frac{e-1}{3})$.\\
\textbf{Final Answer:} $\frac{2\pi(e-1)}{3}$

% ----- CONCEPT 10 -----
\item Find the volume of a solid bounded by the cone $z = \sqrt{3x^2+3y^2}$ and the sphere $x^2+y^2+z^2=4$ (above the xy-plane).\\
\textbf{Solution:}
Bounds: Sphere is $\rho=2$. Cone is $z = \sqrt{3}\sqrt{x^2+y^2} \implies \rho\cos\phi = \sqrt{3}\rho\sin\phi \implies \tan\phi = 1/\sqrt{3} \implies \phi = \pi/6$. 
The solid is bounded between the z-axis and the cone, so $0 \le \phi \le \pi/6$. $0 \le \theta \le 2\pi$.
$V = \int_0^{2\pi} d\theta \int_0^{\pi/6} \sin\phi d\phi \int_0^2 \rho^2 d\rho$
$V = (2\pi) [-\cos\phi]_0^{\pi/6} [\frac{\rho^3}{3}]_0^2 = 2\pi (1 - \frac{\sqrt{3}}{2}) (\frac{8}{3}) = \frac{16\pi(2-\sqrt{3})}{6} = \frac{8\pi(2-\sqrt{3})}{3}$.\\
\textbf{Final Answer:} $\frac{8\pi(2-\sqrt{3})}{3}$

\item Evaluate $\iiint_B \cos\left((x^2+y^2+z^2)^{3/2}\right) dV$ where $B$ is the ball of radius $\sqrt[3]{\pi/2}$ centered at the origin.\\
\textbf{Solution:}
Bounds: $0 \le \rho \le \sqrt[3]{\pi/2}, 0 \le \phi \le \pi, 0 \le \theta \le 2\pi$.
Integrand: $\cos(\rho^3)$. Multiply by $dV$.
$I = \int_0^{2\pi} d\theta \int_0^\pi \sin\phi d\phi \int_0^{\sqrt[3]{\pi/2}} \cos(\rho^3) \rho^2 d\rho$
1. $\int_0^{2\pi} d\theta = 2\pi$.
2. $\int_0^\pi \sin\phi d\phi = 2$.
3. $\int_0^{\sqrt[3]{\pi/2}} \cos(\rho^3)\rho^2 d\rho$. Let $u=\rho^3, du=3\rho^2 d\rho$. Bounds $0 \to \pi/2$. 
$\int_0^{\pi/2} \frac{1}{3}\cos(u)du = [\frac{1}{3}\sin(u)]_0^{\pi/2} = \frac{1}{3}(1 - 0) = \frac{1}{3}$.
Multiply: $(2\pi)(2)(1/3) = \frac{4\pi}{3}$.\\
\textbf{Final Answer:} $\frac{4\pi}{3}$

\item Evaluate $\iiint_E z^2 dV$, where $E$ is bounded between the spheres $\rho=2, \rho=3$ and the cones $\phi=\pi/4, \phi=\pi/3$.\\
\textbf{Solution:}
Integrand $z^2 = \rho^2\cos^2\phi$. Multiply by $dV$.
$I = \int_0^{2\pi} d\theta \int_{\pi/4}^{\pi/3} (\cos^2\phi)(\sin\phi) d\phi \int_2^3 \rho^4 d\rho$
1. $\int_0^{2\pi} d\theta = 2\pi$.
2. $\int_2^3 \rho^4 d\rho = \frac{243 - 32}{5} = \frac{211}{5}$.
3. Let $u=\cos\phi, du=-\sin\phi d\phi$. Bounds $\cos(\pi/4)=\sqrt{2}/2$ to $\cos(\pi/3)=1/2$. 
$\int_{1/2}^{\sqrt{2}/2} u^2 du = [\frac{u^3}{3}]_{1/2}^{\sqrt{2}/2} = \frac{(\sqrt{2}/2)^3 - (1/2)^3}{3} = \frac{2\sqrt{2}/8 - 1/8}{3} = \frac{2\sqrt{2}-1}{24}$.
Multiply: $(2\pi)(\frac{211}{5})(\frac{2\sqrt{2}-1}{24}) = \frac{211\pi(2\sqrt{2}-1)}{60}$.\\
\textbf{Final Answer:} $\frac{211\pi(2\sqrt{2}-1)}{60}$

% ----- CONCEPT 12 -----
\item Find the z-coordinate of the center of mass ($\bar{z}$) of the solid bounded above by the sphere $x^2+y^2+z^2=z$ and below by the cone $z=\sqrt{x^2+y^2}$. Assume uniform density.\\
\textbf{Solution:}
Sphere: $\rho^2 = \rho\cos\phi \implies \rho=\cos\phi$. Cone: $\phi=\pi/4$.
Bounds: $0 \le \theta \le 2\pi, 0 \le \phi \le \pi/4, 0 \le \rho \le \cos\phi$.
Volume $V = \int_0^{2\pi} d\theta \int_0^{\pi/4} \sin\phi d\phi \int_0^{\cos\phi} \rho^2 d\rho = 2\pi \int_0^{\pi/4} \frac{\cos^3\phi}{3}\sin\phi d\phi$.
Let $u=\cos\phi$. $V = \frac{2\pi}{3} \int_{\sqrt{2}/2}^1 u^3 du = \frac{2\pi}{3} [\frac{1}{4} - \frac{1}{16}] = \frac{2\pi}{3} (\frac{3}{16}) = \frac{\pi}{8}$.
Z-moment $M_{xy} = \iiint z dV = \int_0^{2\pi} d\theta \int_0^{\pi/4} \int_0^{\cos\phi} (\rho\cos\phi)(\rho^2\sin\phi) d\rho d\phi$.
$M_{xy} = 2\pi \int_0^{\pi/4} \cos\phi\sin\phi [\frac{\rho^4}{4}]_0^{\cos\phi} d\phi = \frac{\pi}{2} \int_0^{\pi/4} \cos^5\phi\sin\phi d\phi$.
Let $u=\cos\phi$. $M_{xy} = \frac{\pi}{2} \int_{\sqrt{2}/2}^1 u^5 du = \frac{\pi}{2} [\frac{u^6}{6}]_{\sqrt{2}/2}^1 = \frac{\pi}{12} [1 - (\frac{\sqrt{2}}{2})^6] = \frac{\pi}{12} [1 - \frac{1}{8}] = \frac{\pi}{12} (\frac{7}{8}) = \frac{7\pi}{96}$.
$\bar{z} = \frac{M_{xy}}{V} = \frac{7\pi/96}{\pi/8} = \frac{7}{12}$.\\
\textbf{Final Answer:} $\bar{z} = \frac{7}{12}$

\item Find the z-coordinate of the center of mass of the uniform solid bounded above by the sphere $\rho=R$ and below by the cone $\phi=\pi/3$.\\
\textbf{Solution:}
Bounds: $0 \le \theta \le 2\pi, 0 \le \phi \le \pi/3, 0 \le \rho \le R$.
$V = \int_0^{2\pi} d\theta \int_0^{\pi/3} \sin\phi d\phi \int_0^R \rho^2 d\rho = 2\pi [-\cos\phi]_0^{\pi/3} [\frac{R^3}{3}] = \frac{2\pi R^3}{3} (1 - \frac{1}{2}) = \frac{\pi R^3}{3}$.
$M_{xy} = \int_0^{2\pi} d\theta \int_0^{\pi/3} \cos\phi\sin\phi d\phi \int_0^R \rho^3 d\rho = 2\pi [\frac{\sin^2\phi}{2}]_0^{\pi/3} [\frac{R^4}{4}] = \frac{\pi R^4}{2} (\frac{(\sqrt{3}/2)^2}{2}) = \frac{\pi R^4}{2} (\frac{3/4}{2}) = \frac{3\pi R^4}{16}$. Wait, $[\frac{\sin^2\phi}{2}]_0^{\pi/3} = \frac{3/4}{2} = 3/8$. So $2\pi (3/8) (R^4/4) = \frac{6\pi R^4}{32} = \frac{3\pi R^4}{16}$.
$\bar{z} = \frac{M_{xy}}{V} = \frac{3\pi R^4/16}{\pi R^3/3} = \frac{9R}{16}$.\\
\textbf{Final Answer:} $\bar{z} = \frac{9R}{16}$

% ----- CONCEPT 13 -----
\item \textbf{Find the Flaw:} To calculate the integral of $f(x,y,z) = z$ over the upper unit hemisphere, a student sets up the integral as:
$\int_0^{2\pi} \int_0^{\pi} \int_0^1 (\rho\cos\phi) \rho^2\sin\phi \,d\rho \,d\phi \,d\theta$.
Identify the flaw and provide the correct step.\\
\textbf{Solution:} 
\textit{Critique:} The phrase "upper hemisphere" implies $z \ge 0$. The angle $\phi$ measures downwards from the positive z-axis. The upper hemisphere restricts $\phi$ to range from $0$ to $\pi/2$. The bound $\pi$ given in the student's setup maps out the entire sphere, not just the upper half. 
\textit{Correct Step:} Change the $\phi$ upper bound to $\pi/2$.\\
\textbf{Final Answer:} The $\phi$ limits should be $0 \le \phi \le \pi/2$.

\item \textbf{Find the Flaw:} A student is asked to find the volume of the region bounded above by the sphere $\rho=4$ and below by the flat horizontal plane $z=2$. The student sets the bounds as $2 \le \rho \le 4$, $0 \le \phi \le \pi/3$, and $0 \le \theta \le 2\pi$. Identify the mathematical flaw in these bounds.\\
\textbf{Solution:} 
\textit{Critique:} The lower bound is a flat plane $z=2$. In spherical coordinates, $z = \rho\cos\phi$. Setting $z=2$ gives $\rho\cos\phi = 2$, which means $\rho = 2\sec\phi$. The student incorrectly assumed that a flat horizontal lower bound translates to a constant radial bound $\rho=2$, which actually represents a smaller sphere, not a flat plane. 
\textit{Correct Step:} The lower limit for $\rho$ must be $2\sec\phi$.\\
\textbf{Final Answer:} The $\rho$ bounds should be $2\sec\phi \le \rho \le 4$.

% ----- CONCEPT 11 -----
\item \textbf{Application (Asset Shock):} In a 3D quantitative risk model, the probability density function (PDF) of a financial shock is spherically symmetric and given by $f(\rho) = C e^{-\rho^3}$ defined over all 3D space ($0 \le \rho < \infty$). Find the normalization constant $C$ such that the total probability integrates to 1.\\
\textbf{Solution:}
We must satisfy $\iiint_{\mathbb{R}^3} f(\rho) dV = 1$.
$1 = \int_0^{2\pi} d\theta \int_0^\pi \sin\phi d\phi \int_0^\infty C e^{-\rho^3} \rho^2 d\rho$
$1 = (2\pi) (2) (C) \int_0^\infty \rho^2 e^{-\rho^3} d\rho$
Let $u=\rho^3, du=3\rho^2 d\rho$.
$1 = 4\pi C \int_0^\infty \frac{1}{3} e^{-u} du = \frac{4\pi C}{3} [-e^{-u}]_0^\infty = \frac{4\pi C}{3} (0 - (-1)) = \frac{4\pi C}{3}$.
Solve for C: $C = \frac{3}{4\pi}$.\\
\textbf{Final Answer:} $C = \frac{3}{4\pi}$

\item \textbf{Application (Expected Distance):} Using the normalized PDF from the previous problem, $f(\rho) = \frac{3}{4\pi} e^{-\rho^3}$, find the expected Euclidean distance $E[\rho]$ of a shock from the origin. (Recall $E[\rho] = \iiint \rho \cdot f(\rho) dV$).\\
\textbf{Solution:}
$E[\rho] = \int_0^{2\pi} d\theta \int_0^\pi \sin\phi d\phi \int_0^\infty \rho \left(\frac{3}{4\pi} e^{-\rho^3}\right) \rho^2 d\rho$
$E[\rho] = (4\pi) \frac{3}{4\pi} \int_0^\infty \rho^3 e^{-\rho^3} d\rho = 3 \int_0^\infty \rho^3 e^{-\rho^3} d\rho$.
Let $u=\rho^3, du=3\rho^2 d\rho \implies \rho = u^{1/3}$.
$E[\rho] = 3 \int_0^\infty u e^{-u} (\frac{du}{3\rho^2}) = 3 \int_0^\infty u e^{-u} \frac{du}{3u^{2/3}} = \int_0^\infty u^{1/3} e^{-u} du$.
This is exactly the definition of the Gamma function $\Gamma(4/3)$.\\
\textbf{Final Answer:} $\Gamma(4/3)$

% ----- SYNTHESIS -----
\item \textbf{Synthesis:} Find the volume of the "spherical cap" bounded above by the sphere $\rho=4$ and below by the plane $z=2$. Set up and evaluate the integral purely in spherical coordinates.\\
\textbf{Solution:}
Intersection: $z=2 \implies \rho\cos\phi=2$. On the sphere $\rho=4$, $4\cos\phi=2 \implies \cos\phi=1/2 \implies \phi=\pi/3$.
Bounds: $\theta \in [0, 2\pi]$. $\phi$ sweeps from $0$ (the z-axis) down to $\pi/3$ (the intersection ring). For a fixed $\phi$, $\rho$ goes from the flat plane ($2\sec\phi$) up to the sphere (4).
$V = \int_0^{2\pi} d\theta \int_0^{\pi/3} \int_{2\sec\phi}^4 \rho^2\sin\phi \,d\rho \,d\phi$
$V = 2\pi \int_0^{\pi/3} [\frac{\rho^3}{3}]_{2\sec\phi}^4 \sin\phi \,d\phi = \frac{2\pi}{3} \int_0^{\pi/3} (64 - 8\sec^3\phi)\sin\phi \,d\phi$.
Break into two parts:
1. $\int_0^{\pi/3} 64\sin\phi \,d\phi = 64[-\cos\phi]_0^{\pi/3} = 64(-\frac{1}{2} - (-1)) = 32$.
2. $\int_0^{\pi/3} 8\sec^3\phi\sin\phi \,d\phi = 8 \int_0^{\pi/3} \frac{\sin\phi}{\cos^3\phi} d\phi$. Let $u=\cos\phi, du=-\sin\phi d\phi$.
$8 \int_{1/2}^1 u^{-3} du = 8 [\frac{u^{-2}}{-2}]_{1/2}^1 = -4 [1 - 4] = 12$.
Combine: $V = \frac{2\pi}{3} [32 - 12] = \frac{40\pi}{3}$.\\
\textbf{Final Answer:} $\frac{40\pi}{3}$

\end{enumerate}

\vspace{1cm}
\hrule \vspace{0.5cm}

\section*{Concept Checklist (End)}
Below is the list of core concepts mapped to the corresponding problems in this set:
\begin{itemize}
    \item \textbf{C1:} Point Conversion (Spherical to Cartesian) $\implies$ \textbf{Problems 1, 2}
    \item \textbf{C2:} Point Conversion (Cartesian to Spherical) $\implies$ \textbf{Problems 3, 4}
    \item \textbf{C3:} Identifying Surfaces from Spherical Equations $\implies$ \textbf{Problems 5, 6}
    \item \textbf{C4:} Converting Cartesian Equations to Spherical Equations $\implies$ \textbf{Problems 7, 8}
    \item \textbf{C5:} Geometric Visualization \& Setting Up Inequalities for 3D Regions $\implies$ \textbf{Problems 9, 10}
    \item \textbf{C6:} Evaluating Separable Triple Integrals (Constant Bounds / Spherical Rectangles) $\implies$ \textbf{Problems 11, 12, 13}
    \item \textbf{C7:} Evaluating Triple Integrals over Spheres and Balls $\implies$ \textbf{Problems 15, 16, 17}
    \item \textbf{C8:} Evaluating Triple Integrals bounded by Cones and Spheres $\implies$ \textbf{Problems 14, 21, 23, 30}
    \item \textbf{C9:} Converting Cartesian Triple Integrals to Spherical Integrals $\implies$ \textbf{Problems 18, 19, 20}
    \item \textbf{C10:} Advanced Integration Techniques (Integration by Parts, Odd Powers of Sine/Cosine) $\implies$ \textbf{Problems 13, 20, 22}
    \item \textbf{C11:} Real-World Applications \& Modeling (Probability \& Expected Distance) $\implies$ \textbf{Problems 28, 29}
    \item \textbf{C12:} Center of Mass / Centroids $\implies$ \textbf{Problems 24, 25}
    \item \textbf{C13:} "Find the Flaw" / Diagnostic Testing $\implies$ \textbf{Problems 26, 27}
\end{itemize}

\end{document}